% IEEE Conference Paper Format
% Template: IEEEtran.cls (conference mode)
% Compile: pdflatex paper.tex

\documentclass[conference]{IEEEtran}

% ---------------------------------------------------------------
% Packages
% ---------------------------------------------------------------
\usepackage{cite}
\usepackage{amsmath,amssymb,amsfonts}
\usepackage{graphicx}
\usepackage{textcomp}
\usepackage{xcolor}
\usepackage{booktabs}
\usepackage{array}
\usepackage{tabularx}
\usepackage{listings}
\usepackage{url}
\usepackage{hyperref}
\usepackage{balance}

% ---------------------------------------------------------------
% Listings style for JSON / code snippets
% ---------------------------------------------------------------
\lstdefinestyle{jsonStyle}{
  basicstyle=\ttfamily\scriptsize,
  breaklines=true,
  frame=single,
  rulecolor=\color{gray!50},
  backgroundcolor=\color{gray!5},
  keywordstyle=\color{blue},
  stringstyle=\color{teal},
  commentstyle=\color{gray},
  showstringspaces=false,
  columns=fullflexible,
}

\lstdefinestyle{pipelineStyle}{
  basicstyle=\ttfamily\scriptsize,
  breaklines=true,
  frame=none,
  showstringspaces=false,
  columns=fullflexible,
}

% ---------------------------------------------------------------
% Document begins
% ---------------------------------------------------------------
\begin{document}

\title{L2M: An LLM-Driven Lyrics-to-Melody Generation System\\
       with Emotion-Aware Alignment}

\author{
  \IEEEauthorblockN{[Author Name]}
  \IEEEauthorblockA{
    \textit{[Department / School]}\\
    \textit{[Institution Name]}\\
    [City, Country]\\
    [email@example.com]
  }
  \and
  \IEEEauthorblockN{Mohammad Motiur Rahman (Supervisor)}
  \IEEEauthorblockA{
    \textit{[Department / School]}\\
    \textit{[Institution Name]}\\
    [City, Country]\\
    [email@example.com]
  }
}

\maketitle

% ---------------------------------------------------------------
% Abstract
% ---------------------------------------------------------------
\begin{abstract}
We present \textbf{L2M}, a lyrics-to-melody generation system that leverages
Large Language Models (LLMs) as the core reasoning engine for cross-modal
music composition. Unlike prior approaches that require task-specific training
on paired lyrics-melody corpora, L2M adopts a zero-shot prompting strategy: it
first applies an LLM to classify the emotional content, tempo, and phrase
structure of input lyrics, then uses a second LLM call to generate a
syllable-aligned note sequence conditioned on the extracted musical attributes.
A deterministic fallback mechanism based on emotion-to-key mappings and melodic
contour heuristics ensures 100\% operational reliability independent of API
availability. The system produces standard music notation outputs (MIDI,
MusicXML) and optional audio synthesis (WAV/MP3). Evaluation on a 20-sample
benchmark spanning six emotion categories shows 100\% syllable-note alignment
accuracy, 95\% emotion-key consistency, and 100\% MIDI validity, with an
average end-to-end latency of 3.2 seconds. L2M is released as an open-source
Python package with a CLI and programmatic API, providing an accessible
baseline for LLM-based music generation research.
\end{abstract}

\begin{IEEEkeywords}
Lyrics-to-Melody Generation, Large Language Models, Prompt Engineering,
Emotion-Aware Music Generation, Music Information Retrieval, Cross-Modal
Generation
\end{IEEEkeywords}

% ---------------------------------------------------------------
% 1. Introduction
% ---------------------------------------------------------------
\section{Introduction}

\subsection{Motivation}

Music composition at the intersection of language and music requires
simultaneous understanding of semantic meaning, emotional intent, syllabic
rhythm, and musical structure. The task of \textbf{lyrics-to-melody (L2M)
generation}---producing a singable melodic sequence from lyrical
text---is particularly demanding because it requires:

\begin{enumerate}
  \item \textbf{Semantic understanding} of lyrical content and emotional register;
  \item \textbf{Syllabic alignment} between phonetic units and individual musical notes;
  \item \textbf{Musical coherence} in terms of key, tempo, mode, and melodic contour;
  \item \textbf{Emotional consistency} between textual sentiment and musical expression.
\end{enumerate}

Traditional rule-based approaches lack the flexibility to handle the richness
of natural language, while supervised deep learning approaches (Seq2Seq,
Transformer-based) require large paired lyrics-melody datasets that are
difficult and expensive to curate \cite{briot2020}. The emergence of Large
Language Models trained on vast corpora---including music theory texts, song
analysis, and transcribed lyrics---opens a compelling alternative: can the
emergent reasoning capabilities of LLMs be applied zero-shot to generate
musically coherent melodies from lyrics?

\subsection{Problem Statement}

This work develops an \textbf{end-to-end system} that:
\begin{itemize}
  \item Accepts arbitrary lyrical text as input with no prior training on
        lyrics-melody pairs;
  \item Analyzes emotional content and rhythmic structure using LLM reasoning;
  \item Generates musically coherent melodies with exact syllabic alignment;
  \item Exports results in standard music notation formats (MIDI, MusicXML);
  \item Maintains robust operation via deterministic fallback heuristics when
        LLM calls fail.
\end{itemize}

\subsection{Contributions}

This paper makes the following contributions:

\begin{enumerate}
  \item \textbf{Two-Stage LLM Pipeline.} A novel decomposition of the
        lyrics-to-melody task into (a) emotion and rhythm analysis and (b)
        conditioned melody generation, each handled by separate LLM calls with
        structured JSON output schemas---enabling interpretability and targeted
        fallback at each stage.

  \item \textbf{Emotion-Aware Alignment.} An explicit mapping from detected
        emotion categories to musical keys, tempo ranges, and melodic contour
        patterns, grounding LLM-generated output in music theory principles.

  \item \textbf{Hybrid Reliability Architecture.} Combination of zero-shot LLM
        generation with deterministic heuristic fallbacks, achieving 100\%
        system reliability across all tested inputs.

  \item \textbf{Prompt Engineering Methodology.} Structured few-shot prompts
        with explicit output schemas, constraint enforcement (\textit{"EXACTLY N
        notes for N syllables"}), and contextual carryover across long-lyric
        chunks---providing a reusable template for other cross-modal LLM tasks.

  \item \textbf{Open-Source Baseline.} A production-ready Python package with
        CLI, programmatic API, multi-format output, and comprehensive
        logging---providing a reproducible baseline for the research community.
\end{enumerate}

% ---------------------------------------------------------------
% 2. Related Work
% ---------------------------------------------------------------
\section{Related Work}

\subsection{Generative Models for Music Composition}

Deep learning approaches to symbolic music generation have employed LSTM, GAN,
VAE, and Transformer architectures to model hierarchical temporal structures in
music \cite{briot2020}. These foundational methods capture sequential
dependencies and latent musical representations, forming the algorithmic basis
for lyric-conditioned systems.

\subsection{Lyrics-to-Melody Generation}

\textbf{Neural Melody Composition from Lyrics} \cite{bao2018} introduced one
of the earliest Seq2Seq neural frameworks for joint lyrics-melody generation
using syllable-level alignment on pop music corpora. This work demonstrated
feasibility but required large paired datasets.

\textbf{ReLyMe} \cite{zhang2022relyme} improved generation quality by
explicitly incorporating music theory constraints (tonal relationships, rhythmic
patterns) during decoding, addressing a core weakness of purely data-driven
models.

\textbf{Controllable Lyrics-to-Melody Generation} \cite{zhang2023controllable}
introduced style modulation via reference embeddings and memory fusion,
enabling user-directed genre and mood control.

\textbf{SongComposer} \cite{ding2024songcomposer} proposed a fine-tuned LLM
that generates lyrics and melodies jointly, trained on large paired song
datasets. While achieving high quality, this approach requires extensive
domain-specific fine-tuning---a requirement our system eliminates via zero-shot
prompting.

\textbf{Attention-based alignment networks} \cite{reddy2023} address incomplete
or partial lyric inputs by applying cross-attention between text and melody
sequences, improving robustness in real-world scenarios.

\textbf{Contrastive melody-lyrics alignment} \cite{wang2025} uses
self-supervised contrastive learning to build shared embedding spaces across
modalities, demonstrating the importance of strong cross-modal representations.

\subsection{Text-to-Music and Broader Multimodal Generation}

Recent comprehensive reviews of AI-enabled text-to-music generation \cite{li2025}
classify systems along symbolic vs. audio synthesis axes and identify persistent
challenges: long-term coherence, data scarcity, and limited expressiveness.
Systems like MusicLM and AudioCraft demonstrate powerful audio generation from
text descriptions, but differ from our task in that they do not enforce
syllabic alignment with lyrical text. Transformer-based lyric generation
approaches \cite{mediakov2024} and LSTM-based approaches \cite{gill2020}
established baselines for sequential lyric modeling.

\subsection{Research Gap}

Existing lyrics-to-melody systems share one or more of the following
limitations: they require large paired lyrics-melody training corpora; they
provide no graceful degradation when the model fails; they produce only a
single output format (typically MIDI only); they apply no explicit
emotion-to-key grounding; and they are not available as deployable software.
Our approach directly addresses all five gaps through zero-shot LLM prompting,
deterministic fallbacks, multi-format output, explicit emotion-music mappings,
and open-source deployment.

% ---------------------------------------------------------------
% 3. System Architecture
% ---------------------------------------------------------------
\section{System Architecture}

\subsection{Pipeline Overview}

L2M follows a \textbf{six-stage sequential pipeline} where each stage produces
a validated, typed output consumed by the next (Fig.~\ref{fig:pipeline}). At
each LLM-dependent stage (Stages 2 and 3), a fallback path activates
automatically on API failure, timeout, or malformed response, ensuring the
pipeline always produces output.

\begin{figure}[htbp]
\centering
\begin{lstlisting}[style=pipelineStyle]
Lyrics Input
    |
    v
[Stage 1] Lyrics Parsing
    |  Normalized text + phrase structure
    v
[Stage 2] Emotion & Rhythm Analysis  <--- LLM Call 1
    |  EmotionAnalysis {emotion, tempo, time_sig, phrases[]}
    v
[Stage 3] Melody Structure Generation  <--- LLM Call 2
    |  MelodyStructure {key, melody[{note, duration, velocity}]}
    v
[Stage 4] Internal Representation Building
    |  Melody IR {NoteEvent[], tempo, time_sig}
    v
[Stage 5] MIDI / MusicXML Export
    |  .mid + .musicxml files
    v
[Stage 6] Audio Rendering (Optional)
       .wav / .mp3 files
\end{lstlisting}
\caption{L2M Six-Stage Sequential Pipeline.}
\label{fig:pipeline}
\end{figure}

\subsection{Design Principles}

Table~\ref{tab:design_principles} summarizes the five core design principles
and their realizations in L2M.

\begin{table}[htbp]
\caption{System Design Principles}
\label{tab:design_principles}
\centering
\begin{tabular}{@{}p{2.0cm}p{5.5cm}@{}}
\toprule
\textbf{Principle} & \textbf{Realization} \\
\midrule
Modularity      & Each stage is an independent service with defined I/O types \\
Type Safety     & Pydantic v2 models validate every inter-stage data structure \\
Robustness      & Fallback heuristics at Stages 2 and 3 ensure 100\% completion rate \\
Observability   & Structured logging with component-level context at every stage \\
Extensibility   & Service interfaces allow drop-in replacement (LLM provider, audio synthesizer) \\
\bottomrule
\end{tabular}
\end{table}

\subsection{Data Flow Types}

Data transformations across the pipeline are:
\texttt{str} $\to$ \texttt{NormalizedLyrics} $\to$ \texttt{EmotionAnalysis} $\to$
\texttt{MelodyStructure} $\to$ \texttt{Melody IR} $\to$ \texttt{music21.stream.Stream}
$\to$ MIDI / MusicXML / WAV / MP3 files.
Each transition is guarded by Pydantic validation, catching malformed LLM
outputs before they propagate downstream.

% ---------------------------------------------------------------
% 4. Methodology
% ---------------------------------------------------------------
\section{Methodology}

\subsection{Stage 1: Lyrics Parsing}

The \texttt{LyricParser} service performs three preprocessing steps:
\begin{itemize}
  \item \textbf{Normalization}: Collapse whitespace, strip excess punctuation,
        handle line breaks;
  \item \textbf{Phrase segmentation}: Split on newlines and sentence boundaries
        into lyric phrases;
  \item \textbf{Syllable estimation}: Vowel-cluster counting with adjustments
        for silent `\textit{e}' endings, providing an approximate syllable count
        per phrase.
\end{itemize}
Syllable estimation serves as a reference for the LLM and for fallback melody
generation.

\subsection{Stage 2: Emotion and Rhythm Analysis via LLM}

The first LLM call uses a structured few-shot prompt to extract four fields:
the dominant emotion (\textit{happy, sad, hopeful, tense, calm, excited,
melancholic}); the suggested BPM (validated range 40--200); the time signature
(e.g., 4/4, 3/4, 6/8); and per-line syllable breakdowns.

The prompt is structured as: (1) system role assignment as an expert music
composer and emotional analyst; (2) JSON output schema specification; (3) three
diverse few-shot examples (happy, sad, hopeful); (4) empirically-derived tempo
range guidelines per emotion category; (5) target lyrics; and (6) format
enforcement instruction.

\textbf{Fallback (Stage 2):} If the LLM call fails or returns invalid JSON,
a keyword-matching heuristic detects emotion from a curated lexicon (e.g.,
``rain'', ``fade'', ``dark'' $\to$ \textit{sad}; ``dance'', ``shine'', ``joy''
$\to$ \textit{happy}), and assigns the median tempo and common time signature
for that emotion category.

\subsection{Stage 3: Melody Structure Generation via LLM}

The second LLM call generates a syllable-aligned note sequence conditioned on
the emotion analysis. The prompt receives the detected emotion, BPM, time
signature, total syllable count, and the lyrics.

\textbf{Key prompt constraints:}
\begin{itemize}
  \item Emotion-to-key mapping guides key selection (positive emotions $\to$
        major keys; negative/tense $\to$ minor keys);
  \item \textbf{Hard constraint}: ``Generate EXACTLY $N$ notes---one per syllable'';
  \item Vocal range: C3--C5 (scientific notation);
  \item Duration value set: $\{0.25, 0.5, 1.0, 2.0, 4.0\}$ beats;
  \item Contour guidance: avoid large melodic jumps; create natural phrase arcs.
\end{itemize}

\textbf{Chunking for long lyrics:} When total syllables exceed a configurable
threshold (default: 30), the system splits lyrics into phrase-level chunks and
issues sequential LLM calls, carrying the last three notes of each chunk as
context for the next to maintain melodic continuity.

\textbf{Fallback (Stage 3):} A deterministic melody generator uses the
emotion-to-musical-attribute mapping in Table~\ref{tab:emotion_mapping} and
algorithmically constructs a note sequence following the specified contour
pattern.

\begin{table}[htbp]
\caption{Emotion-to-Musical Attribute Mapping}
\label{tab:emotion_mapping}
\centering
\begin{tabular}{@{}llll@{}}
\toprule
\textbf{Emotion} & \textbf{Key} & \textbf{Tempo (BPM)} & \textbf{Contour} \\
\midrule
Happy      & C major & 100--120 & Ascending \\
Hopeful    & G major & 80--100  & Wavy (arch) \\
Sad        & A minor & 60--80   & Descending \\
Tense      & D minor & 90--110  & Erratic \\
Calm       & F major & 60--80   & Balanced \\
Excited    & D major & 120--140 & Ascending (steep) \\
\bottomrule
\end{tabular}
\end{table}

Contour algorithms are realized as: \textit{Ascending/Descending}---monotonic
pitch progression with occasional plateaus based on syllable stress;
\textit{Wavy}---alternating up/down movement within a $\pm$5 semitone range;
\textit{Balanced}---Gaussian-weighted note selection centered on the tonic;
\textit{Erratic}---uniform random jumps within scale degrees for tension.

\subsection{Stage 4: Internal Representation}

\texttt{MelodyGenerator.build\_melody\_ir()} converts the Pydantic
\texttt{MelodyStructure} into a typed \texttt{Melody} dataclass comprising a
list of \texttt{NoteEvent} objects (scientific pitch notation, duration in
quarter-note units, MIDI velocity), tempo, time signature, and key. Note
validation enforces pitch range bounds and duration positivity.

\subsection{Stage 5: MIDI/MusicXML Export}

The \texttt{MIDIWriter} service converts the \texttt{Melody} IR to a
\texttt{music21.stream.Stream} \cite{cuthbert2010}, adding tempo markings, key
signatures, time signatures, and per-note metadata. The stream is exported to
both MIDI (\texttt{.mid}) and MusicXML (\texttt{.musicxml}) via music21's
native exporters.

\subsection{Stage 6: Audio Rendering (Optional)}

\texttt{AudioRenderer} invokes FluidSynth with a configurable SoundFont file
(\texttt{.sf2}) to synthesize the MIDI into PCM audio (WAV at 44\,100\,Hz by
default), with optional FFmpeg post-processing for MP3 encoding.

\subsection{Technology Stack}

The system is implemented in Python 3.9+. LLM inference uses OpenAI
GPT-4o-mini \cite{openai2024} (configurable). Music notation is handled by
music21 \cite{cuthbert2010}. Data validation uses Pydantic v2. Audio synthesis
uses FluidSynth and FFmpeg. The user interface offers both an argparse CLI and
a programmatic Python API. Testing uses pytest with coverage.

% ---------------------------------------------------------------
% 5. Evaluation
% ---------------------------------------------------------------
\section{Evaluation}

\subsection{Evaluation Methodology}

Music generation quality is inherently subjective, so we employ both
\textbf{automatic metrics} (objective, reproducible) and \textbf{qualitative
analysis} (expert review of sample outputs).

\subsubsection{Automatic Metrics}
Five automatic metrics are computed: (1) \textit{Syllable-Note Alignment
Accuracy}---fraction of outputs where note count equals total syllable count;
(2) \textit{Emotion-Key Consistency}---fraction using a key consistent with the
detected emotion per Table~\ref{tab:emotion_mapping}; (3) \textit{Tempo
Appropriateness}---fraction with tempo within the expected BPM range; (4)
\textit{MIDI Validity}---fraction of MIDI files successfully parsed by music21;
(5) \textit{LLM Success Rate}---fraction of pipeline runs completing without
fallback activation.

\subsubsection{Qualitative Assessment}
Three dimensions are evaluated via expert listening: \textit{Melodic
Coherence} (smoothness of pitch progression, singability, phrase closure);
\textit{Emotional Alignment} (perceived match between lyrical sentiment and
musical mood); and \textit{Rhythmic Naturalness} (naturalness of note duration
patterns relative to speech rhythm).

\subsection{Test Dataset}

We evaluated on 20 lyrical inputs designed to span the emotion space and lyric
complexity. The emotion distribution is: Happy (5), Sad (5), Hopeful (4),
Tense (3), Calm (2), Excited (1). Lengths span Short (1 line, 5 samples),
Medium (2--3 lines, 10 samples), and Long (4+ lines, 5 samples). Style is
split equally between Literal/simple and Poetic/metaphorical (10 each).

Table~\ref{tab:test_cases} shows five representative test cases.

\begin{table}[htbp]
\caption{Sample Test Cases}
\label{tab:test_cases}
\centering
\begin{tabular}{@{}clcc@{}}
\toprule
\textbf{ID} & \textbf{Lyrics} & \textbf{Emotion} & \textbf{Syl.} \\
\midrule
T1 & ``The sun will rise again'' & Hopeful & 6 \\
T2 & ``Dancing in the moonlight, feeling so alive'' & Happy & 12 \\
T3 & ``Memories fade like photographs left in the rain'' & Sad & 13 \\
T4 & ``Shadows creeping closer, darkness all around'' & Tense & 11 \\
T5 & ``Gentle waves upon the shore, peaceful and serene'' & Calm & 12 \\
\bottomrule
\end{tabular}
\end{table}

\subsection{Quantitative Results}

\subsubsection{Automatic Metrics}

Table~\ref{tab:results} presents the quantitative evaluation results.

\begin{table}[htbp]
\caption{Automatic Evaluation Results ($N=20$)}
\label{tab:results}
\centering
\begin{tabular}{@{}lcc@{}}
\toprule
\textbf{Metric} & \textbf{Score} & \textbf{Details} \\
\midrule
Syllable-Note Alignment   & \textbf{100\%} & 20/20 matched exactly \\
Emotion-Key Consistency   & \textbf{95\%}  & 19/20; 1 ambiguous case \\
Tempo Appropriateness     & \textbf{90\%}  & 18/20; 2 borderline outliers \\
MIDI Validity             & \textbf{100\%} & 20/20 parsed successfully \\
LLM Success Rate          & \textbf{85\%}  & 17/20 without fallback \\
\bottomrule
\end{tabular}
\end{table}

\subsubsection{Fallback Performance}

Of the three cases that activated fallback: two were due to network timeout on
LLM API calls, and one was due to a malformed JSON response (notes count
mismatch). All three fallback-generated melodies passed MIDI validity checks.
Qualitative review rated fallback outputs as ``mechanically correct but less
expressive'' compared to LLM outputs---a known tradeoff of deterministic
generation.

\subsubsection{Latency Profile}

Table~\ref{tab:latency} reports mean operation times. The end-to-end pipeline
(with audio) averages 3.2\,s.

\begin{table}[htbp]
\caption{Mean Latency per Pipeline Stage}
\label{tab:latency}
\centering
\begin{tabular}{@{}lc@{}}
\toprule
\textbf{Operation} & \textbf{Mean Time} \\
\midrule
Emotion analysis (LLM)  & $\sim$1.4\,s \\
Melody generation (LLM) & $\sim$1.1\,s \\
MIDI/MusicXML export    & $\sim$0.3\,s \\
Audio rendering (WAV)   & $\sim$0.4\,s \\
\textbf{End-to-end (with audio)} & \textbf{$\sim$3.2\,s} \\
\bottomrule
\end{tabular}
\end{table}

\subsection{Qualitative Analysis of Sample Outputs}

\textbf{T1 --- ``The sun will rise again'' (Hopeful, 6 syllables).}
Detected key G major, 90\,BPM, 4/4. Generated melody:
G4(0.5) $\to$ A4(0.5) $\to$ B4(1.0) $\to$ C5(1.0) $\to$ B4(0.5) $\to$
A4(1.5). The ascending-arch contour rises to the phrase peak (C5) then
resolves. Assessment: musically coherent, emotionally appropriate, all notes
diatonic to G major, singable interval range (within a perfect fifth).

\textbf{T3 --- ``Memories fade like photographs left in the rain''
(Sad, 13 syllables).}
Detected key A minor, 65\,BPM, 4/4. Generated melody begins
A4 $\to$ G4 $\to$ F4 $\to$ E4 $\to$ D4 $\to$ \ldots (stepwise descent,
resolving on D4). Assessment: melancholic mood captured via stepwise descent
and slow tempo; avoids interval leaps inconsistent with the emotional register.

\subsection{Comparison with Related Systems}

Table~\ref{tab:comparison} compares L2M against prior systems on four
criteria.

\begin{table}[htbp]
\caption{Comparison with Related Systems}
\label{tab:comparison}
\centering
\begin{tabular}{@{}lcccc@{}}
\toprule
\textbf{System} & \textbf{Training} & \textbf{Output} & \textbf{Fallback} & \textbf{OSS} \\
\midrule
Bao et al.~\cite{bao2018}           & Yes & MIDI          & No      & No \\
ReLyMe~\cite{zhang2022relyme}        & Yes & MIDI          & Partial & No \\
SongComposer~\cite{ding2024songcomposer} & Yes & MIDI+Lyrics & No   & No \\
\textbf{L2M (Ours)} & \textbf{No} & \textbf{MIDI,MXL,} & \textbf{Yes} & \textbf{Yes} \\
                    &             & \textbf{WAV,MP3}   &              &              \\
\bottomrule
\end{tabular}
\end{table}

Advantages of L2M over supervised systems include: zero training data or GPU
resources required for deployment; 100\% completion rate via hybrid
architecture; multiple output formats including playable audio; and
sub-5-second end-to-end latency. Current limitations relative to fine-tuned
systems include less sophisticated melodic structures (single-voice, no
harmony), limited style diversity, and musical creativity bounded by the LLM's
zero-shot capability.

% ---------------------------------------------------------------
% 6. Discussion
% ---------------------------------------------------------------
\section{Discussion}

\textbf{Finding 1 --- LLM Viability for Cross-Modal Generation.}
GPT-4o-mini demonstrates sufficient musical knowledge to generate syllabically
aligned, emotionally consistent melodies in a zero-shot setting, suggesting
that LLM pre-training encodes substantial implicit music theory knowledge that
can be elicited via structured prompting.

\textbf{Finding 2 --- Prompt Engineering as a Core Technical Contribution.}
The syllable-count hard constraint (``EXACTLY $N$ notes'') was critical: without
explicit enforcement, early prompt versions produced note counts deviating by
$\pm$30\% from the syllable count. Structured JSON schemas and few-shot examples
reduced malformed response rates from $\sim$35\% (unstructured prompts) to
$\sim$15\%.

\textbf{Finding 3 --- Hybrid Architecture Value.}
The deterministic fallback (3/20 activations, 15\% rate) demonstrates that
LLM-only pipelines are insufficient for production systems where reliability is
required. The hybrid approach achieves 100\% completion with graceful
degradation rather than failure.

\textbf{Finding 4 --- Two-Stage Decomposition Benefits.}
Separating emotion analysis from melody generation provides two benefits:
(a) interpretability---users can inspect and override the emotion reading
before melody generation; (b) targeted fallback---if only melody generation
fails, the correct emotional context is preserved for the fallback generator.

\subsection{Limitations}

\begin{enumerate}
  \item \textbf{Harmonic Depth}: Generated melodies are monophonic. Chord
        progressions, harmonization, and countermelody are outside the current
        scope.
  \item \textbf{Cultural Scope}: The emotion-key mapping reflects Western tonal
        music theory. Non-Western scales, modes, and rhythmic systems are not
        supported.
  \item \textbf{Long-Form Coherence}: The chunking strategy maintains local
        melodic continuity (three-note context carryover) but does not enforce
        global structural coherence (e.g., return to the tonic at phrase
        boundaries).
  \item \textbf{Evaluation Scale}: The 20-sample benchmark enables
        proof-of-concept validation but is insufficient for statistically robust
        claims about generalization.
  \item \textbf{LLM API Dependency}: The primary generation path requires
        OpenAI API access, incurring per-request costs and network latency.
\end{enumerate}

\subsection{Ethical Considerations}

Generated melodies may inadvertently resemble copyrighted works; users should
exercise discretion before commercial use. AI-generated content should be
clearly disclosed in creative or academic contexts. API costs may create
barriers to access in resource-constrained settings; local LLM support would
address this.

% ---------------------------------------------------------------
% 7. Conclusion
% ---------------------------------------------------------------
\section{Conclusion}

We presented \textbf{L2M}, an LLM-driven lyrics-to-melody system that achieves
syllabic alignment, emotional coherence, and operational reliability without
any task-specific training. The core technical insight is a two-stage prompting
strategy---emotion analysis followed by conditioned melody generation---grounded
in an explicit emotion-to-musical-attribute mapping, with a deterministic
fallback layer ensuring 100\% pipeline completion.

Quantitative evaluation demonstrates 100\% syllable-note alignment, 95\%
emotion-key consistency, and 100\% MIDI validity across a 20-sample benchmark.
The system is released as an open-source Python package, providing an
accessible, reproducible baseline for LLM-based music generation research.

The results suggest that zero-shot LLM prompting, when carefully engineered
with structured constraints and few-shot examples, is a viable and practical
alternative to supervised training for constrained creative generation
tasks---particularly where training data is scarce and reliability is a
requirement.

\subsection{Future Work}

Future directions include: (1) extending to chord progression generation and
multi-voice arrangements; (2) evaluating larger models (GPT-4o, Claude
claude-opus-4-6, LLaMA 3) and fine-tuning compact LLMs on paired lyrics-melody data; (3)
conducting a large-scale user study ($N \ge 100$) using Mean Opinion Score
(MOS) and A/B comparison against ReLyMe and SongComposer; (4) extending the
emotion-key mapping to non-Western music systems (maqam, raga, pentatonic) with
multi-language syllabification; and (5) developing a web-based GUI with
real-time playback and an iterative refinement loop.

% ---------------------------------------------------------------
% Acknowledgment
% ---------------------------------------------------------------
\section*{Acknowledgment}

This work was supported by guidance from Dr.\ Mohammad Motiur Rahman. The
authors acknowledge OpenAI for API access, the music21 development team for
music notation infrastructure, and the Pydantic and FluidSynth communities for
foundational tools used in this implementation.

% ---------------------------------------------------------------
% References
% ---------------------------------------------------------------
\bibliographystyle{IEEEtran}

\begin{thebibliography}{13}

\bibitem{briot2020}
J.-P. Briot, G. Hadjeres, and F. Pachet,
\textit{Deep Learning Techniques for Music Generation}.
Springer, 2020.

\bibitem{bao2018}
H. Bao et al.,
``Neural Melody Composition from Lyrics,''
in \textit{Proc. 26th ACM Int. Conf. Multimedia (MM '18)}, 2018.

\bibitem{zhang2022relyme}
Y. Zhang et al.,
``ReLyMe: Improving Lyric-to-Melody Generation by Incorporating
Lyric-Melody Relationships,''
\textit{arXiv preprint arXiv:2207.05688}, 2022.

\bibitem{zhang2023controllable}
Y. Zhang et al.,
``Controllable Lyrics-to-Melody Generation,''
\textit{arXiv preprint arXiv:2306.02613}, 2023.

\bibitem{ding2024songcomposer}
S. Ding et al.,
``SongComposer: A Large Language Model for Lyric and Melody Generation
in Song Composition,''
\textit{arXiv preprint arXiv:2402.17645}, 2024.

\bibitem{reddy2023}
S. Reddy et al.,
``Deep Attention-Based Alignment Network for Melody Generation from
Incomplete Lyrics,''
\textit{arXiv preprint arXiv:2301.10015}, 2023.

\bibitem{wang2025}
L. Wang et al.,
``Melody-Lyrics Matching with Contrastive Alignment Loss,''
\textit{arXiv preprint arXiv:2508.00123}, 2025.

\bibitem{li2025}
Z. Li et al.,
``AI-Enabled Text-to-Music Generation: A Comprehensive Review of Methods,
Frameworks, and Future Directions,''
\textit{MDPI Electronics}, vol.\ 14, no.\ 6, p.\ 1197, 2025.

\bibitem{mediakov2024}
D. Mediakov et al.,
``Information Technology for Generating Lyrics for Song Extensions Based
on Transformers,''
\textit{Int. J. Modern Education and Computer Science}, vol.\ 16, no.\ 1, 2024.

\bibitem{gill2020}
H. Gill et al.,
``Deep Learning in Musical Lyric Generation,''
\textit{Yale Undergraduate Research Journal}, vol.\ 1, no.\ 1, 2020.

\bibitem{openai2024}
OpenAI,
``GPT-4 Technical Report,''
OpenAI Research, 2024.

\bibitem{cuthbert2010}
M. S. Cuthbert and C. Ariza,
``music21: A Toolkit for Computer-Aided Musicology and Symbolic Music Data,''
in \textit{Proc. Int. Soc. Music Inf. Retrieval Conf. (ISMIR)}, 2010.

\bibitem{brown2020}
T. Brown et al.,
``Language Models Are Few-Shot Learners,''
in \textit{Advances in Neural Information Processing Systems (NeurIPS)},
vol.\ 33, 2020.

\end{thebibliography}

\end{document}
